\SGASection{2}{Existence theorem: construction of a piece of group}
\SGARefTarget{sga3:xxv:section:2}

\SGANextSectionTarget{sga3:xxv:subsection:2-1}\SGASubsection{2.1}{}
\SGARefTarget{sga3:xxv:paragraph:2:1}

Choose once and for all a splitting of \(G_k\), denoted
\[
  (G_k,T_k,M,R)
\]
(see \SGARefLink{sga3:xxii:def:1.13}{Exposé XXII, 1.13}), a system \(\Delta\) of simple roots of \(R\)
(defining the system \(R^+\) of positive roots), and a Chevalley system
\((X_{\alpha,k})_{\alpha\in R}\) of \(G_k\) (\SGARefLink{sga3:XXIII:6.1}{Exposé XXIII, 6.1} and \SGARefLink{sga3:XXIII:6.2}{6.2})
satisfying the additional condition (see \SGARefLink{sga3:xx:2.6}{Exposé XX, 2.6})
\[
  X_{\alpha,k}X_{-\alpha,k}=1\qquad(\alpha\in R).
\]
Finally, choose on the subgroup of \(M\) generated by \(R\) a total order
compatible with the group structure and such that the roots \(>0\) are the
elements of \(R^+\). We then write the roots as
\[
 -\alpha_n<-\alpha_{n-1}<\cdots<-\alpha_1
 <\alpha_1<\alpha_2<\cdots<\alpha_n.
\]
For \(\alpha\in R\), let \(U_{\alpha,k}\) denote the vector group
corresponding to the root \(\alpha\), and let
\[
 p_{\alpha,k}\colon\mathbb G_{a,k}\xrightarrow{\ \sim\ }U_{\alpha,k}
\]
be the isomorphism of vector groups defined by \(X_{\alpha,k}\).

\SGANextSectionTarget{sga3:xxv:subsection:2-2}\SGASubsection{2.2}{}
\SGARefTarget{sga3:xxv:paragraph:2:2}

The splitting of \(G_k\) includes, in particular, an isomorphism of
\(k\)-groups
\[
 T_k\simeq D_k(M).
\]
Put \(T=D(M)\). This is a \(\mathbb Z\)-torus, and the preceding isomorphism
may be regarded as an isomorphism
\[
 T_k\simeq T\otimes_{\mathbb Z}k.
\]
We have
\[
 \operatorname{Hom}_{\mathbb Z\text{-gr.}}(T,\mathbb G_m)=M,
\]
and shall regard the elements of \(R\subset M\) as characters of \(T\).
Likewise, the elements of \(R^\ast\) will be regarded as morphisms of
\(\mathbb Z\)-groups \(\mathbb G_m\to T\).

\SGANextSectionTarget{sga3:xxv:subsection:2-3}\SGASubsection{2.3}{}
\SGARefTarget{sga3:xxv:paragraph:2:3}

For each \(\alpha\in R^+\), let \(\mathbb G_a(\alpha)\) be a copy of
\(\mathbb G_a\), and consider the \(\mathbb Z\)-scheme
\[
 U=\mathbb G_a(\alpha_1)\times\cdots\times\mathbb G_a(\alpha_n).
\]
If \(U_k\) denotes the unipotent part of the Borel subgroup \(B_k\) of
\(G_k\) defined by \(R^+\), let
\[
 a\colon U\otimes_{\mathbb Z}k\xrightarrow{\ \sim\ }U_k
\]
be the isomorphism of \(k\)-schemes defined by
\[
 a(x_1,\ldots,x_n)
 =p_{\alpha_1,k}(x_1)\cdots p_{\alpha_n,k}(x_n).
\]

\SGANextSectionTarget{sga3:xxv:subsection:2-4}\SGASubsection{2.4}{}
\SGARefTarget{sga3:xxv:paragraph:2:4}

The group law of \(U_k\) is expressed by relations of the form
\[
 a(x_1,\ldots,x_n)a(y_1,\ldots,y_n)=a(z_1,\ldots,z_n),
\]
where, for \(h=1,\ldots,n\), each \(z_h\) is a polynomial
\[
 z_h=x_h+y_h+
 Q_h(x_1,\ldots,x_{h-1},y_1,\ldots,y_{h-1}),
\]
whose coefficients are integers (\SGARefLink{sga3:xxii:lem:5.5.8}{Exposé XXII, 5.5.8}, and \SGARefLink{sga3:XXIII:6.4}{Exposé XXIII,
6.4}). Moreover,
\[
 Q_h(x_1,\ldots,x_{h-1},0,\ldots,0)=0.
\]

Equip \(U\) with the composition law defined by these formulas, which are
indeed ``defined over \(\mathbb Z\)''. Since this law induces the group law
on \(U_k\), it is associative, and \((0)\) is a unit element. Indeed, these
two assertions are expressed by relations among the polynomials \(Q_h\),
and \(\mathbb Z\to k\) is injective. This is a group law: if \((x_i)\) is
a section of \(U\) over an arbitrary scheme \(S\), its inverse \((y_i)\)
is computed recursively by
\[
 y_i=-x_i-Q_i(x_1,\ldots,x_{i-1},y_1,\ldots,y_{i-1}),
\]
formulas which are again ``defined over \(\mathbb Z\)''.

In summary, we have constructed a group law on \(U\) for which the
isomorphism \(a\) above is an isomorphism of groups.

For each \(\alpha\in R^+\), consider the morphism
\[
 p_\alpha\colon\mathbb G_a\longrightarrow U
\]
defined by \(p_\alpha(x)=(x_i)\), where
\[
 x_i=
 \begin{cases}
 x,&\alpha_i=\alpha,\\
 0,&\alpha_i\ne\alpha.
 \end{cases}
\]
It is a closed immersion and a group homomorphism; denote its image by
\(U_\alpha\). Since
\[
 (x_i)=p_{\alpha_1}(x_1)\cdots p_{\alpha_n}(x_n),
\]
\(U\) is identified with the product
\[
 U=U_{\alpha_1}U_{\alpha_2}\cdots U_{\alpha_n}.
\]

\SGANextSectionTarget{sga3:xxv:subsection:2-5}\SGASubsection{2.5}{}
\SGARefTarget{sga3:xxv:paragraph:2:5}

Let \(T=D(M)\) act on each \(U_\alpha\) through the character \(\alpha\).
One verifies immediately that this defines an action of \(T\) on the
group \(U\), and one may form the semidirect product \(B=T\ltimes U\).
There is a canonical isomorphism of \(k\)-groups
\[
 B\otimes_{\mathbb Z}k\xrightarrow{\ \sim\ }B_k.
\]

If an arbitrary order is now placed on \(R^+\), the morphism
\[
 \prod_{\alpha\in R^+}U_\alpha\longrightarrow U
\]
defined by multiplication in \(U\) is still an isomorphism. Indeed, both
sides are flat \(\mathbb Z\)-schemes of finite presentation, so it is
enough to verify the assertion on geometric fibers. We are then reduced
to Lazard's theory (\textit{Bible}, \S13.1, \SGARefLink{sga3:xxv:paragraph:2:1}{Proposition 1}): regard \(U\)
as a group with operators \(T\), and use the fact that the \(U_\alpha\)
are pairwise nonisomorphic as groups with operators, since the characters
\(\alpha\in R^+\) of \(T\) are pairwise distinct on every fiber.

